% !TeX root = ../main.tex

% ==================================================
% 2.12. PyTorch nn.Module의 객체 모델 — 지금까지 배운 Python 클래스 개념의 통합
% ==================================================

2.2부터 2.11까지 Python의 클래스와 객체 모델을 학습하였다.

이번 2.12는 새로운 문법을 많이 배우는 단계라기보다, 지금까지 학습한 개념들이 실제 딥러닝 프레임워크에서 어떻게 결합되는지 확인하는 단계다.

특히 다음 코드를 단순히 암기하는 것이 아니라, 각 문장이 어떤 객체를 생성하고, 어디에 저장하며, 어떤 탐색과 호출 과정을 거치는지 설명할 수 있게 되는 것이 목표다.

\begin{dhkcode}
import torch
import torch.nn as nn


class MyModel(nn.Module):
    def __init__(self):
        super().__init__()

        self.linear1 = nn.Linear(10, 20)
        self.relu = nn.ReLU()
        self.linear2 = nn.Linear(20, 5)

    def forward(self, x):
        x = self.linear1(x)
        x = self.relu(x)
        x = self.linear2(x)
        return x


model = MyModel()
x = torch.randn(2, 10)
y = model(x)
\end{dhkcode}

이 코드에서 우리가 추적할 것은 크게 네 가지다.

\begin{dhksmry}
1. 클래스와 인스턴스가 어떻게 생성되는가?

2. self.linear1 같은 하위 모듈은 어디에 저장되는가?

3. model(x)는 어떻게 forward()로 연결되는가?

4. PyTorch는 파라미터와 하위 모듈을 어떻게 관리하는가?
\end{dhksmry}

\subsubsection{\zsh{nn.Module}은 무엇인가?}

먼저 \zsh{nn.Module} 자체를 살펴보자.

\begin{dhkcode}
import torch.nn as nn

print(type(nn.Module))
\end{dhkcode}

실행 결과:

\begin{dhkzsh}
<class 'type'>
\end{dhkzsh}

\zsh{nn.Module}은 클래스 객체다.

따라서 다음 정의는

\begin{dhkcode}
class MyModel(nn.Module):
    pass
\end{dhkcode}

Python의 일반적인 상속과 동일한 원리로 동작한다.

\begin{dhkcode}
print(MyModel.__bases__)
print(MyModel.__mro__)
\end{dhkcode}

출력 예시:

\begin{dhkzsh}
(<class 'torch.nn.modules.module.Module'>,)

(<class '__main__.MyModel'>,
 <class 'torch.nn.modules.module.Module'>,
 <class 'object'>)
\end{dhkzsh}

즉, \zsh{MyModel}은 \zsh{nn.Module}을 상속하는 새로운 클래스 객체다.

\zsh{nn.Module}이 가진 메서드와 특수 메서드들은 \zsh{MyModel}의 MRO를 통해 사용될 수 있다.

예를 들어 다음 메서드들이 있다.

\begin{dhkcode}
model.parameters()
model.state_dict()
model.train()
model.eval()
model.to("cpu")
\end{dhkcode}

이 메서드들을 우리가 \zsh{MyModel} 안에 직접 정의하지 않았더라도 사용할 수 잇는 이유는 상속 때문이다.

\begin{dhknote}
\begin{dhksmry}[centered]
nn.Module의 핵심 역할
\end{dhksmry}
\end{dhknote}

\zsh{nn.Module}은 단순히 \zsh{forward()}를 호출하는 기능만 제공하지 않는다.

모델 내부의 상태와 구조를 관리하는 기반 클래스다.

대표적으로 다음을 관리한다.

\begin{dhktable}
|관리 대상|	내부 등록 위치|	예시|
|:---|:---|:---|
|하위 모듈|	_modules|	nn.Linear, nn.ReLU|
|학습 파라미터|	_parameters|	nn.Parameter|
|버퍼|	_buffers|	BatchNorm의 running statistics|
\end{dhktable}

이 세 가지는 앞으로 모델을 이해할 때 계속 등장하게 될 것이다.

\subsubsection{\zsh{model = MyModel()} - 객체 생성 과정}

먼저 다음 코드만 살펴보자.

\begin{dhkcode}
class MyModel(nn.Module):
    def __init__(self):
        super().__init__()
        self.linear = nn.Linear(10, 5)


model = MyModel()
\end{dhkcode}

2.6에서 학습한 내용을 적용하면 \zsh{MyModel()}은 클래스 객체를 호출하는 표현식이다.

일반적인 클래스 생성 과정은 다음과 같다.

\begin{dhkgraph}[uniform]
node n1 (1,1) :: MyModel() // 클래스 객체 호출
node n2 (2,1) :: type.__call__ // 일반적인 클래스 생성 메커니즘
node n3 (3,1) :: __new__ // MyModel 인스턴스 생성
node n4 (4,1) [accent] :: MyModel.__init__ // 인스턴스 초기화
node n5 (5,1) :: model // 생성된 인스턴스에 이름 바인딩

edge n1 -> n2
edge n2 -> n3
edge n3 -> n4
edge n4 -> n5
\end{dhkgraph}

여기서 \zsh{model}이라는 이름은 새로 생성된 \zsh{MyModel} 인스턴스를 참조하게 된다.

\begin{dhkcode}
print(type(model) is MyModel)
print(isinstance(model, nn.Module))
\end{dhkcode}

실행 결과:

\begin{dhkzsh}
True
True
\end{dhkzsh}

이것은 일반적인 Python 객체 생성과 다르지 않다.

차이는 초기화 과정에서 \zsh{nn.Module}이 제공하는 내부 관리 구조가 설정된다는 점이다.

\subsubsection{\zsh{super().__init__()}은 왜 먼저 실행해야 하는가?}

다음 코드를 보자.

\begin{dhkcode}
class MyModel(nn.Module):
    def __init__(self):
        super().__init__()
        self.linear = nn.Linear(10, 5)
\end{dhkcode}

2.9에서 배운 것처럼 \zsh{super()}는 MRO를 기준으로 다음 구현을 탐색한다.

이 경우 MRO는 다음과 같다.

\begin{dhksmry}[centered]
MyModel \ra nn.Module \ra object
\end{dhksmry}

따라서 \zsh{super().__init__()}은 \zsh{nn.Module.__init__()}으로 연결된다.

이 초기화 과정에서 PyTorch가 모듈 등록과 상태 관리에 필요한 내부 구조를 준비한다.

대표적으로 다음과 같은 속성들이 있다.

\begin{dhkcode}
self._parameters
self._buffers
self._modules
\end{dhkcode}

이들은 보통 딕셔너리 형태의 내부 등록 공간이다.

초기화 이후 다음처럼 확인할 수 있다.

\begin{dhkcode}
class MyModel(nn.Module):
    def __init__(self):
        super().__init__()


model = MyModel()

print(type(model._parameters))
print(type(model._buffers))
print(type(model._modules))
\end{dhkcode}

일반적인 실행 결과:

\begin{dhkzsh}
<class 'dict'>
<class 'dict'>
<class 'dict'>
\end{dhkzsh}

\begin{dhknote}
\begin{dhksmry}[centered]
초기화 순서를 바꾸면?
\end{dhksmry}
\end{dhknote}

다음 코드를 보자.

\begin{dhkcode}
class BadModel(nn.Module):
    def __init__(self):
        self.linear = nn.Linear(10, 5)
        super().__init__()

model = BadModel()
\end{dhkcode}

일반적인 PyTorch에서는 다음과 같은 오류가 발생한다.

\begin{dhkzsh}
AttributeError: cannot assign module before Module.__init__() call
\end{dhkzsh}

왜 이런 문제가 발생할까?

\zsh{self.linear}에 \zsh{nn.Linear} 인스턴스를 대입하는 작업은 단순한 Python 속성 대입이 아니기 때문이다.

\zsh{nn.Module}은 하위 모듈을 자동으로 등록하는 메커니즘을 제공한다.

그런데 그 등록을 위한 내부 상태가 아직 준비되지 않았다면 하위 모듈을 등록할 수 없다.

따라서 일반적인 PyTorch 모델에서는 다음 순서를 지켜야 한다.

\begin{dhkcode}
def __init__(self):
    super().__init__()

    self.linear = nn.Linear(10, 5)
\end{dhkcode}

\zsh{super().__init__()}은 단순한 관례가 아니라, PyTorch 모듈의 등록 시스템을 사용하기 위한 초기화 단계다.

\subsubsection{\zsh{self.lienar = nn.Linear(...)}는 일반적인 속성 대입인가?}

이번 학습에서 가장 중요한 부분 중 하나다.

다음 코드를 보자.

\begin{dhkcode}
self.linear = nn.Linear(10, 5)
\end{dhkcode}

일반적인 Python 클래스라면 다음과 같은 동작을 기대할 수 있다.

\begin{dhkcode}
self.__dict__["linear"] = ...
\end{dhkcode}

하지만 \zsh{nn.Module}에서는 조금 다르다.

2.7에서 학습했던 \zsh{__setattr__()}이 개입하기 때문이다.

\zsh{nn.Module}은 \zsh{__setattr__()}을 재정의해서 대입되는 객체의 종류에 따라 등록 방식을 다르게 처리한다.

예를 들어 \zsh{nn.Module} 인스턴스를 속성에 대입하면 하위 모듈로 등록할 수 있다.

실제로 확인해 보자.

\begin{dhkcode}
import torch.nn as nn


class MyModel(nn.Module):
    def __init__(self):
        super().__init__()

        self.linear = nn.Linear(10, 5)


model = MyModel()
\end{dhkcode}

이제 인스턴스 딕셔너리를 확인해 보자.

\begin{dhkcode}
print("linear" in model.__dict__)
print("linear" in model._modules)
\end{dhkcode}

실행 결과:

\begin{dhkzsh}
False
True
\end{dhkzsh}

\zsh{linear}라는 이름은 일반적인 인스턴스 딕셔너리에 직접 저장되지 않았다.

대신 \zsh{_modules}에 등록되어 있다.

\begin{dhkcode}
print(model._modules)
\end{dhkcode}

출력 예시:

\begin{dhkzsh}
{'linear': Linear(in_features=10, out_features=5, bias=True)}
\end{dhkzsh}

이 결과는 중요한 의미를 가진다.

\zsh{self.linear}이라는 속성 접근 문법을 사용한다고 해서 반드시 \zsh{linear}라는 이름이 \zsh{self.__dict__}에 저장되어 있는 것은 아니다.

2.7에서 학습한 일반적인 속성 탐색 구조가 PyTorch에서는 사용자 정의 속성 접근 메커니즘에 의해 확장된 것이다.

\subsubsection{\zsh{_modules}에 저장된 객체는 어떻게 \zsh{self.linear}로 접근되는가?}

여기서 자연스럽게 의문이 생긴다.

\zsh{linear}가 \zsh{model.__dict__}에 없다면 다음 코드는 어떻게 동작할까?

\begin{dhkcode}
model.linear
\end{dhkcode}

그 이유는 \zsh{nn.Module}이 \zsh{__getattr__()}도 제공하기 때문이다.

일반적인 속성 탐색으로 원하는 속성을 찾지 못하면 \zsh{__getattr__()}이 호출될 수 있다고 배웠다.

PyTorch는 이 메커니즘을 활용해 등록된 파라미터, 버퍼, 화위 모듈을 속성처럼 접근할 수 있게 한다.

개념적인 흐름은 다음과 같다.

\begin{dhkgraph}[uniform]
node n1 (1,1) :: model.linear // 속성 접근
node n2 (2,1) :: 일반 속성 탐색 // 클래스/인스턴스 속성 확인
node n3 (3,1) :: nn.Module.__getattr__ // 일반 탐색에서 찾지 못한 경우
node n4 (4,1) :: model._modules['linear'] // 등록된 하위 모듈 조회
node n5 (5,1) :: nn.Linear 인스턴스 // 객체 반환

edge n1 -> n2
edge n2 -> n3
edge n3 -> n4
edge n4 -> n5

caption :: 일반적인 등록 하위 모듈에 대한 속성 접근을 단순화한 흐름
\end{dhkgraph}

즉, 다음 두 표현식은 일반적으로 같은 하위 모듈 객체를 참조한다.

\begin{dhkcode}
print(model.linear is model._modules["linear"])
\end{dhkcode}

실행 결과:

\begin{dhkzsh}
True
\end{dhkzsh}

여기서 우리가 학습한 두 특수 메서드가 연결된다.

\begin{dhksmry}
    __setattr__() : 하위 모듈 등록

    __getattr__() : 등록된 하위 모듈 조회
\end{dhksmry}

이것이 PyTorch가 일반적인 Python 속성 문법을 사용하면서도 내부적으로 모델 구조를 관리할 수 있는 이유 중 하나다.

\subsubsection{\zsh{nn.Linear} 인스턴스의 파라미터는 어디에 저장되는가?}

이제 다음 코드로 넘어가 보자.

\begin{dhkcode}
linear = nn.Linear(10,5)
\end{dhkcode}

\zsh{linear}는 \zsh{nn.Linear} 클래스의 인스턴스다.

이 인스턴스에는 대표적으로 다음 파라미터가 있다.

\begin{dhkcode}
linear.weight
linear.bias
\end{dhkcode}

각각의 타입을 확인해 보자.

\begin{dhkcode}
print(type(linear.weight))
print(type(linear.bias))
\end{dhkcode}

실행 결과:

\begin{dhkzsh}
<class 'torch.nn.parameter.Parameter'>
<class 'torch.nn.parameter.Parameter'>
\end{dhkzsh}

여기서 \zsh{nn.parameter}는 일반적인 텐서와 관련된 특수한 클래스다.

\zsh{torch.Tensor}를 상속하며, \zsh{nn.Module}의 속성으로 대입되면 파라미터로 자동 등록될 수 있다.

\begin{dhknote}
\begin{dhksmry}[centered]
내부 등록 위치 확인
\end{dhksmry}
\end{dhknote}

\begin{dhkcode}
print("weight" in linear.__dict__)
print("weight" in linear._parameters)
\end{dhkcode}

실행 결과:

\begin{dhkzsh}
False
True
\end{dhkzsh}

즉, \zsh{weight}는 \zsh{_parameters}에 등록되어 있다.

\begin{dhkcode}
print(linear._parameters.keys())
\end{dhkcode}

실행 결과:

\begin{dhkzsh}
dict_keys(['weight', 'bias'])
\end{dhkzsh}

여기서 중요한 점은 모든 텐서가 자동으로 파라미터로 등록되는 것은 아니라는 것이다.

다음 코드를 비교해 보자.

\begin{dhkcode}
class Example(nn.Module):
    def __init__(self):
        super().__init__()

        self.a = torch.randn(3)
        self.b = nn.Parameter(torch.randn(3))

example = Example()

print("a" in example._parameters)
print("b" in example._parameters)
\end{dhkcode}

실행 결과:

\begin{dhkzsh}
False
True
\end{dhkzsh}

일반 텐서 \zsh{a}는 자동으로 학습 파라미터로 등록되지 않는다.

반면 \zsh{nn.parameter} 객체인 \zsh{b}는 파라미터로 등록된다.

이 차이는 오빝마이저와 모델 상태 관리에서 매우 중요하다.

\subsubsection{\zsh{_modules, _parameters, _buffers}의 차이}

지금까지 두 종류의 등록 공간을 확인했다.

\begin{dhksmry}
    _modules : 하위 모듈

    _parameters : 파라미터
\end{dhksmry}

여기에 하나가 더 있다.

바로 \zsh{_buffers}다.

버퍼는 모델 상태의 일부이지만 일반적으로 옵티마이저가 학습하는 파라미터는 아닌 텐서를 관리할 때 사용한다.

예를 들어 Batch Normalization의 running mean과 running variance가 대표적이다.

직접 버퍼를 등록할 수도 있다.

\begin{dhkcode}
class Example(nn.Module):
    def __init__(self):
        super().__init__()

        self.linear = nn.Linear(3, 2)

        self.scale = nn.Parameter(
            torch.ones(2)
        )

        self.register_buffer(
            "running_value",
            torch.zeros(2)
        )

model = Example()
\end{dhkcode}

이제 내부 등록 구조를 확인해 보자.

\begin{dhkcode}
print(model._modules.keys())
print(model._parameters.keys())
print(model._buffers.keys())
\end{dhkcode}

실행 결과:

\begin{dhkzsh}
dict_keys(['linear'])
dict_keys(['scale'])
dict_keys(['running_value'])
\end{dhkzsh}

이 구조를 정리하면 다음과 같다.

\begin{dhkgraph}[uniform]
node n1 (1,2) :: nn.Module 인스턴스
node n2 (2,1) :: linear // 하위 모듈
node n3 (2,2) :: scale // 학습 파라미터
node n4 (2,3) :: running_value // 비파라미터 상태
group g1(n1) [accent] :: model
group g2(n2) :: _modules
group g3(n3) :: _parameters 
group g4(n4) :: _buffers
split g1 -> g2, g3, g4

caption :: 하위 모듈인 linear 역시 자체 등록 공간을 가진다. // linear._parameters -> weight, bias
\end{dhkgraph}

여기서 중요한 구조가 하나 더 있다.

\zsh{model.linear}는 또 다른 \zsh{nn.Module} 인스턴스다.

따라서 \zsh{model.linear}도 자체적인 \zsh{_parameters, _modules, _buffers}를 가질 수 있다.

이것이 PyTorch 모델이 계층적으로 구성되는 기반이다.

\subsubsection{\zsh{model.parameters()}는 어떻게 하위 모듈의 파라미터를 찾는가?}

다음 모델을 다시 보자.

\begin{dhkcode}
class MyModel(nn.Module):
    def __init__(self):
        super().__init__()

        self.linear1 = nn.Linear(10, 20)
        self.relu = nn.ReLU()
        self.linear2 = nn.Linear(20, 5)

    def forward(self, x):
        x = self.linear1(x)
        x = self.relu(x)
        x = self.linear2(x)
        return x

model = MyModel()
\end{dhkcode}

이제 다음 코드를 실행해 보자.

\begin{dhkcode}
for name, parameter in model.named_parameters():
    print(name, tuple(parameter.shape))
\end{dhkcode}

실행 결과:

\begin{dhkzsh}
linear1.weight (20, 10)
linear1.bias (20,)
linear2.weight (5, 20)
linear2.bias (5,)
\end{dhkzsh}

여기서 주목할 점은 \zsh{MyModel} 자체의 \zsh{_parameters}에는 직접 등록한 파라미터가 없을 수 있다는 것이다.

\begin{dhkcode}
print(model._parameters)
\end{dhkcode}

실행 결과:

\begin{dhkzsh}
{}
\end{dhkzsh}

그런데도 \zsh{model.named_parameters()}는 네 개의 파라미터를 찾는다.

왜 그럴까?

\zsh{nn.Module}이 등록된 하위 모듈들을 재귀적으로 탐색할 수 있기 때문이다.

개념적으로 다음 구조를 탐색하는 것이다.

\begin{dhkdir}[in=6]
model
>linear1
>>weight
>>bias
>relu
>linear2
>>weight
>>bias
\end{dhkdir}

\zsh{relu}에는 일반적으로 학습 파라미터가 없으므로 \zsh{named_parameters()} 결과에 나타나지 않는다.

\begin{dhknote}
\begin{dhksmry}[centered]
중요한 구분
\end{dhksmry}
\end{dhknote}

\begin{dhkcode}
model.parameters()
\end{dhkcode}

는 파라미터를 순회할 수 있는 Iterator를 반환한다.

\begin{dhkcode}
model.named_parameters()
\end{dhkcode}

는 파라미터 이름과 파라미터 객체의 쌍을 순회할 수 있는 Iterator를 반환한다.

따라서 다음 코드를 사용할 수 있다.

\begin{dhkcode}
optimizer = torch.optim.Adam(
    model.parameters(),
    lr=0.001
)
\end{dhkcode}

dㅗㅂ티마이저는 이 Iterator를 통해 모델의 등록된 파라미터를 전달받는 것이다.

여기서 2.10에서 학습한 Iterator 프로토콜이 다시 등장한다.

\subsubsection{\zsh{state_dict()}는 무엇을 저장하는가?}

다음 코드를 보자.

\begin{dhkcode}
state = model.state_dict()
\end{dhkcode}

\zsh{state_dict()}는 모델의 상태를 이름과 텐서 값의 대응 관계로 제공한다.

기본적으로 등록되 ㄴ파라미터와 영속적(persistent) 버퍼가 포함된다.

예를 들어 앞의 \zsh{MyModel}에서는 다음 키를 얻을 수 있다.

\begin{dhkcode}
print(list(model.state_dict().keys()))
\end{dhkcode}

실행 결과:

\begin{dhkzsh}
['linear1.weight',
 'linear1.bias',
 'linear2.weight',
 'linear2.bias']
\end{dhkzsh}

여기서 중요한 사실은 다음과 같다.

\zsh{state_dict()}는 모델 객체 전체를 그대로 저장하는 것이 아니라, 모델의 등록된 상태를 구조화해서 제공하는 메커니즘이다.

따라서 모델 구조를 정의하는 클래스와 그 클래스의 파라미터 값은 개념적으로 구분해야 한다.

예를 들어 다음 코드는 모델 상태를 파일로 저장하는 일반적인 방식이다.

\begin{dhkcode}
torch.save(
    model.state_dict(),
    "model.pt"
)
\end{dhkcode}

다시 불러올 때는 먼저 같은 구조의 모델 인스턴스를 생성한 뒤 상태를 로드할 수 있다.

\begin{dhkcode}
model = MyModel()

state = torch.load(
    "model.pt",
    map_location="cpu",
    weights_only=True
)

model.load_state_dict(state)
\end{dhkcode}

여기서 \zsh{MyModel()}은 모델 구조를 생성하고,

\zsh{load_state_dict()}는 등록된 상태를 해당 모델에 반영한다.

이것이 모델의 구조와 상태를 분리하는 방식이다.

\subsubsection{\zsh{model(x)} - 호출 메커니즘과 \zsh{forward()}}

이제 모델 실행으로 넘어가 보자.

\begin{dhkcode}
x = torch.randn(2, 10)

y = model(x)
\end{dhkcode}

여기서 \zsh{model}은 함수 객체가 아니다.

\zsh{MyModel} 인스턴스다.

그런데도 \zsh{model(x)}라는 호출 문법을 사용할 수 있다.

2.10에서 학습한 \zsh{__call__()} 프로토콜 때문이다.

\zsh{nn.Module}은 모델 인스턴스를 호출 간으하게 만드는 메커니즘을 제공한다.

따라서 다음 결과를 얻을 수 있다.

\begin{dhkcode}
print(callable(model))
\end{dhkcode}

실행 결과:

\begin{dhkzsh}
True
\end{dhkzsh}

\begin{dhknote}
\begin{dhksmry}[centered]
호출 과정
\end{dhksmry}
\end{dhknote}

개념적으로 다음과 같은 흐름이다.

\begin{dhkbranch}[in=8]
node main :: model(x)
arrow main
node main :: nn.Module의 호출 메커니즘
arrow main
node main :: MyModel.forward(x)
branchline main -> n1 :: self.linear1(x)
stop n1
branchline main -> n2 :: self.relu(x)
stop n2
branchline main -> n3 :: self.linear2(x)
stop n3
arrow main
node main :: y 변환
\end{dhkbranch}

여기서 중요한 것은 \zsh{forward()}를 우리가 직접 호출하지 않았다는 점이다.

\zsh{model(x)}가 상속된 호출 메커니즘을 통해 \zsh{forward()}를 실행하도록 연결된 것이다.

\begin{dhknote}
\begin{dhksmry}[centered]
실제 메서드 확인
\end{dhksmry}
\end{dhknote}

\begin{dhkcode}
print(model.forward.__self__ is model)
\end{dhkcode}

실행 결과:

\begin{dhkzsh}
True
\end{dhkzsh}

\zsh{forward()}는 인스턴스 메서드이므로 \zsh{model}에 바인딩된다.

또한 다음처럼 확인할 수 있다.

\begin{dhkcode}
print(model.forward.__func__ is MyModel.forward)
\end{dhkcode}

실행 결과:

\begin{dhkzsh}
True
\end{dhkzsh}

즉, \zsh{forward()} 역시 2.5에서 학습한 일반적인 바운드 메서드의 원리로 동작한다.

\subsubsection{\zsh{forward()}를 오버라이딩했는데 왜 \zsh{nn.Module}이 이를 호출할 수 있는가?}

이 부분은 2.8의 상속과 오버라이딩을 연결하는 중요한 지점이다.

다음 코드를 보자.

\begin{dhkcode}
class MyModel(nn.Module):
    def forward(self, x):
        return x * 2
\end{dhkcode}

\zsh{nn.Module}은 호출 메커니즘을 제공하고, 자식 클래스는 \zsh{forward()}를 오버라이딩한다.

그런데 부모 클래스의 호출 메커니즘은 자식 클래스의 \zsh{forward()}를 실행할 수 있다.

왜 그럴까?

호출 메커니즘이 실제 모델 인스턴스의 \zsh{forward} 속성을 조회하기 때문이다.

이때 \zsh{self}는 실제 \zsh{MyModel} 인스턴스다.

따라서 일반적인 속성 탐색 규칙에 따라 \zsh{MyModel.forward}가 선택된다.

이것을 간단한 일반 Python 클래스로 표현하면 다음과 같다.

\begin{dhkcode}
class Base:
    def __call__(self, x):
        return self.forward(x)

    def forward(self, x):
        raise NotImplementedError


class Child(Base):
    def forward(self, x):
        return x * 2

obj = Child()

print(obj(10))
\end{dhkcode}

실행 결과:

\begin{dhkzsh}
20
\end{dhkzsh}

여기서 \zsh{Child}는 \zsh{__call__()}을 정의하지 않았다.

하지만 \zsh{Base.__call__()}을 상속했다.

그리고 \zsh{Base.__call__()} 내부에서 다음 코드를 실행한다.

\begin{dhkcode}
self.forward(x)
\end{dhkcode}

여기서 \zsh{self}는 \zsh{Child} 인스턴스다.

따라서 \zsh{Child.forward()}가 선택된다.

이 구조가 PyTorch의 모델 호출을 이해하기 위한 핵심적인 Python 객체 모델이다.

단, 실제 \zsh{nn.Module}의 호출 메커니즘에는 Hook 처리 등 추가적인 기능이 포함되어 있다.

\subsubsection{\zsh{model.forward(x)} 대신 \zsh{model(x)}를 사용하는 이유}

다음 두 표현식을 비교해 보자.

\begin{dhkcode}
y1 = model(x)

y2 = model.forward(x)
\end{dhkcode}

두 표현식 모두 \zsh{forward()}의 계산을 실행할 수 있다.

하지만 실행 경로는 같지 않다.

첫 번째는 \zsh{nn.Module}의 호출 메커니즘을 사용한다.

두 번째는 \zsh{forward()}를 직접 호출한다.

PyTorch는 모듈 호출 과정에서 여러 Hook을 처리할 수 있다.

예를 들어 다음과 같은 Hook이 있다.

\begin{dhksmry}
    - Forward pre-hook

    - Forward hook

    - Backward 관련 hook
\end{dhksmry}

\zsh{model.forward(x)}를 직접 호출하면 모듈 호출 메커니즘을 우회하므로 이러한 처리가 정상적으로 적용되지 않을 수 있다.

따라서 일반적인 모델 실행에서는 다음 형태를 사용해야 한다.

\begin{dhkcode}
y = model(x)
\end{dhkcode}

이것은 단순한 스타일 문제가 아니라 객체 호출 프로토콜을 통해 프레임워크의 실행 메커니즘을 사용한다는 의미다.

\subsubsection{모델 구조와 계산 그래프는 같은 것인가?}

여기서 중요한 개념적 구분을 하나 해야 한다.

우리는 지금까지 \zsh{nn.Module}이 하위 모듈ㅇ르 등록하고 관리하는 구조를 살펴보았다.

그런데 이 구조와 자동 미분을 위한 계산 그래프는 동일하지 않다.

\begin{dhknote}
\begin{dhksmry}[centered]
모듈 구조
\end{dhksmry}
\end{dhknote}

\begin{dhkdir}[in=6]
MyModel
>linear1
>relu
>linear2
\end{dhkdir}

이 구조는 Python 객체와 하위 모듈의 등록 관계를 나타낸다.

모델 인스턴스가 어떤 하위 모듈ㅇ르 가지고 있는지 설명하는 구조다.

\begin{dhknote}
\begin{dhksmry}[centered]
계산 그래프
\end{dhksmry}
\end{dhknote}

반면 다음 코드가 실행되면

\begin{dhkcode}
x = torch.randn(2, 10)

y = model(x)
\end{dhkcode}

텐서 연산이 수행된다.

\begin{dhksmry}[centered]
x
\da
Linear 연산
\da
ReLU 연산
\da
Linear 연산
\da
y
\end{dhksmry}

이 연산들 중 자동 미분 추적 대상이 되는 연산들은 역전파를 위한 그래프에 반영된다.

여기서 중요한 차이는 다음과 같다.

\begin{dhktable}
|구분|	모듈 구조|	자동 미분 계산 그래프|
|:---|:---|:---|
|기반|	Python 객체 및 모듈 등록|	텐서 연산|
|주요 목적|	모델 구성 및 상태 관리|	미분 계산|
|생성 시점|	주로 모델 초기화 과정|	추적 대상 연산이 실행될 때|
|관리 대상|	모듈, 파라미터, 버퍼|	연산 관계 및 역전파 정보|
|동적 변화|	모듈 구조 변경 가능|	실행 경로에 따라 달라질 수 있음|
\end{dhktable}

특히 PyTorch의 일반적인 eager 실행에서는 연산이 수행되면서 자동 미분 그래프가 구성된다.

따라서 ㅁ듈 등록 구조가 곧 계산 그래프인 것은 아니다.

이 구분은 이후 Transformer 모델을 분석할 때도 중요하다.

\subsubsection{Transformer에서는 이 구조가 어떻게 확장되는가?}

다음은 단순화한 Transformer 블록이다.

\begin{dhkcode}
class TransformerBlock(nn.Module):
    def __init__(self, d_model):
        super().__init__()

        self.attention = nn.MultiheadAttention(
            embed_dim=d_model,
            num_heads=8,
            batch_first=True
        )

        self.norm = nn.LayerNorm(d_model)

        self.ffn = nn.Sequential(
            nn.Linear(d_model, 4 * d_model),
            nn.ReLU(),
            nn.Linear(4 * d_model, d_model)
        )

    def forward(self, x):
        attn_out, _ = self.attention(x, x, x)

        x = self.norm(x + attn_out)

        x = x + self.ffn(x)

        return x
\end{dhkcode}

이 코드는 지금까지 배운 Python 객체 모델을 그대로 사용한다.

먼저 클래스 객체가 생성된다.

\begin{dhkcode}
TransformerBlock
\end{dhkcode}

그다음 인스턴스를 생성한다.

\begin{dhkcode}
block = TransformerBlock(512)
\end{dhkcode}

이 과정에서 하위 모듈들이 등록된다.

\begin{dhkdir}[in=6]
TransformerBlock 인스턴스
> attention
>> MultiheadAttention
> norm
>> LayerNorm
>ffn
>>Sequential
>>>Linear
>>>ReLU
>>>Linear
\end{dhkdir}

이제 다음 코드를 실행하면

\begin{dhkcode}
x = torch.randn(2, 16, 512)

y = block(x)
\end{dhkcode}

\zsh{nn.Module}의 호출 메커니즘을 통해 \zsh{ransformerBlock.forward()}가 실행된다.

그 내부에서 다시 하위 모듈들을 호출한다.

\begin{dhkcode}
self.attention(x, x, x)
self.norm(x + attn_out)
self.ffn(x)
\end{dhkcode}

각 하위 모듈도 \zsh{nn.Module} 인스턴스이므로 동일한 호출 메커니즘을 사요한다.

즉, Transformer는 특별한 Python 객체 시스템 위에서 동작하는 것이 아니라, 우리가 지금까지 학습한 클래스·상속·메서드·특수 메서드·모듈 등록 구조를 계층적으로 조합한 모델이다.

\subsubsection{지금까지의 객체 모델을 하나의 그림으로 정리하기}

다음 코드를 다시 보자.

\begin{dhkcode}
class MyModel(nn.Module):
    def __init__(self):
        super().__init__()

        self.linear = nn.Linear(10, 5)

    def forward(self, x):
        return self.linear(x)


model = MyModel()
y = model(x)
\end{dhkcode}

이 코드를 객체 관계와 실행 관계로 구분해 보면 다음과 같다.

\begin{dhkgraph}[uniform]
node n1 (1,1) :: MyModel 클래스 객체  // nn.Module 상속, forward 정의
node n2 (2,1) [accent] :: model 인스턴스 // _modules["linear"] -> nn.Linear 인스턴스 // linear._parameters -> weight, bias // 상속된 __call__ 메커니즘
node n3 (3,1) :: nn.Module 호출 처리 // Hook 등 프레임워크 실행 처리
node n4 (4,1) [accent] :: MyModel.forward(x) // self.linear(x) -> 텐서 연산
node n5 (5,1) :: 객체 생성 -> 등록 -> 호출 -> 연산

edge n1 -> n2 :: MyModel()
edge n2 -> n3 :: model(x)
edge n3 -> n4
\end{dhkgraph}

이제 코드 한 줄마다 어떤 객체와 메커니즘이 사용되는지 연결할 수 있다.

\subsubsection{최종 확인 - 지금까지의 Python 클래스 학습과 연결}

2.2부터 2.12까지의 내용을 정리해 보자.

\begin{dhktable}
|단원|	핵심 개념|	PyTorch에서의 대응|
|:---|:---|:---|
|2.2|	클래스도 객체|	MyModel, nn.Module|
|2.3|	클래스와 인스턴스|	model = MyModel()|
|2.4|	인스턴스 속성과 네임스페이스|	self.linear|
|2.5|	메서드와 self|	forward(self, x)|
|2.6|	__new__, __init__|	모델 생성 및 초기화|
|2.7|	클래스·인스턴스 속성|	_modules, _parameters, 속성 등록|
|2.8|	상속과 오버라이딩|	nn.Module 상속, forward() 재정의|
|2.9|	super()와 MRO|	super().__init__()|
|2.10|	특수 메서드|	model(x)|
|2.11|	메서드 종류|	forward, from_config, 검증 메서드|
|2.12|	객체 모델 통합|	모델 등록·상태 관리·실행|
\end{dhktable}

\begin{dhknote}
\begin{dhksmry}[centered]
마지막으로 반드시 구분해야 할 다섯 가지
\end{dhksmry}
\end{dhknote}

\begin{dhknote}
\begin{dhksmry}[centered]
첫째, 클래스 객체와 인스턴스 객체는 다르다.
\end{dhksmry}
\end{dhknote}

\begin{dhkcode}
MyModel
model
\end{dhkcode}

\zsh{MyModel}은 클래스 객체이고, \zsh{model}은 그 인스턴스다.

\begin{dhknote}
\begin{dhksmry}[centered]
둘째, 속성 접근 문법과 실제 저장 위치는 다를 수 있다.
\end{dhksmry}
\end{dhknote}

\begin{dhkcode}
model.linear
\end{dhkcode}

는 일반적으로 \zsh{_modules}에 등록된 하위 모듈을 조회한다.

\begin{dhknote}
\begin{dhksmry}[centered]
셋째, 파라미터와 일반 텐서는 등록 방식이 다르다.
\end{dhksmry}
\end{dhknote}

\begin{dhkcode}
self.weight = nn.Parameter(...)
self.tensor = torch.randn(...)
\end{dhkcode}

첫 번째는 학습 파라미터로 등록되고, 두 번째는 일반적으로 자동 등록되지 않는다.

\begin{dhknote}
\begin{dhksmry}[centered]
넷째, model(x)와 model.forward(x)는 실행 경로가 다르다.
\end{dhksmry}
\end{dhknote}

\zsh{model(x)}는 \zsh{nn.Module}의 호출 메커니즘을 사용한다.

\begin{dhknote}
\begin{dhksmry}[centered]
다섯째, 모듈 구조와 계산 그래프는 서로 다른 개념이다.
\end{dhksmry}
\end{dhknote}

모듈 구조는 객체의 구성 관계이고, 자동 미분 계산 그래프는 텐서 연산의 관계이다.

\subsubsection{Python 클래스와 객체 모델 학습 완료}

이번 2.12까지 학습했다면 다음 코드의 의미를 객체 모델 관점에서 설명할 수 있는 기반이 갖춰진 것이다.

\begin{dhkcode}
class MyModel(nn.Module):
    def __init__(self):
        super().__init__()
        self.linear = nn.Linear(10, 5)

    def forward(self, x):
        return self.linear(x)
\end{dhkcode}

핵심은 다음과 같다.

\begin{dhksmry}[centered]
Class Object
\da
Instance Creation
\da
Module Initialization
\da
Attribute Registration
\da
Callable Protocol
\da
Overridden forward
\da
Tensor Operations
\end{dhksmry}

이제 Python 클래스는 단순히 \zsh{class}와 \zsh{self}를 사용하는 문법이 아니라, 객체 생성·네임스페이스·디스크립터·상속·MRO·특수 메서드가 결합된 실행 시스템으로 이해할 수 있다.
